\section{应用 Agent vs 模型能力}

前面讲 PMF，本章回到 AI 应用创业的核心竞争问题：模型越来越强，应用公司还剩什么壁垒？Lovart 的答案是垂直工作流。模型红利会快速同质化，但设计师场景里的上下文、审美、资产管理、版本迭代和交付流程，仍然需要产品来承接。

\begin{figure}[H]
\centering
\includegraphics[width=0.94\textwidth]{agent-app-vs-model.png}
\caption{应用 Agent vs 模型能力：应用创业要在模型红利窗口内沉淀工作流和心智。自制概念图，依据 00:55:00--01:10:30 对谈内容整理。}
\end{figure}

\begin{knowledgebox}{读图：模型红利不是应用壁垒}
底座模型变强会打开窗口，但也会让同类功能快速普及。应用公司要把窗口期转成工作流、用户心智、数据回流和垂直体验。
\end{knowledgebox}

\subsection{通用 Agent 与垂直 Agent 会并存}

上一节说明应用壁垒来自工作流，本节讨论它和通用入口的关系。访谈里提到，通用 Agent 和垂直 Agent 可能并存。通用 Agent 会成为入口，但垂直 Agent 在专业场景中有更深工具、更强上下文和更好体验。Lovart 可能既被通用 Agent 调用，也可能成为设计场景中的专门入口。

\begin{importantbox}{垂直 Agent 的判断标准}
如果一个场景有专业上下文、专门工具、反复迭代、交付标准和愿意付费的用户，就值得做垂直 Agent。设计正符合这些条件。
\end{importantbox}

\begin{knowledgebox}{通用/垂直共存模型}
\begin{tabularx}{\textwidth}{p{0.22\textwidth}p{0.34\textwidth}X}
\toprule
形态 & 角色 & Lovart 的位置 \\
\midrule
通用 Agent 入口 & 接收用户意图，分发任务 & 可能调用 Lovart 完成设计任务。 \\
垂直 Agent 工作台 & 深入专业场景和工具链 & 直接服务设计师和创意团队。 \\
模型底座 & 提供生成、理解和推理能力 & Lovart 需要利用但不能只依赖底座。 \\
插件/工具生态 & 连接不同专业能力 & Lovart 可以成为设计能力节点。 \\
\bottomrule
\end{tabularx}
\end{knowledgebox}

\subsection{商业模式与版权/品牌风险}

通用/垂直关系说清楚后，本节进入商业化和风险。设计类 Agent 的商业模式并不简单。它可能按个人订阅收费，也可能做团队工作台，还可能按项目、素材量或商业授权收费。但越进入真实设计工作，版权、素材来源、品牌一致性和客户资产安全就越重要。设计师和企业客户不会只问“能不能生成”，还会问“生成结果能不能商用”“能不能符合品牌规范”“能不能被团队继续编辑”。

\begin{knowledgebox}{设计 Agent 商业模式表}
\begin{tabularx}{\textwidth}{p{0.20\textwidth}p{0.34\textwidth}X}
\toprule
模式 & 优势 & 风险 \\
\midrule
个人订阅 & 上手快，增长简单 & 客单价低，容易被通用工具替代。 \\
团队工作台 & 进入协作流程，留存更强 & 需要权限、版本、资产管理。 \\
按项目收费 & 和客户价值更相关 & 计价复杂，交付责任更重。 \\
商业授权/素材 & 可以绑定品牌和版权价值 & 法务、版权和责任边界更复杂。 \\
\bottomrule
\end{tabularx}
\end{knowledgebox}

\begin{importantbox}{设计 Agent 的企业级门槛}
真正进入企业设计流程时，Lovart 需要解决的不只是生成质量，还有资产安全、版权可解释、品牌一致、团队协作和可编辑交付。这些才是垂直 Agent 的深水区。
\end{importantbox}

\subsection{应用壁垒的三层结构}

上一节说明商业化会把 Lovart 推向更深工作流，本节拆解它可能形成的壁垒。Lovart 如果要长期抵御模型公司下探，需要三层壁垒。第一层是体验壁垒：让设计师觉得它比通用聊天框更顺手。第二层是工作流壁垒：团队项目、资产、版本和交付都在里面。第三层是生态壁垒：与设计工具、素材库、品牌系统和客户协作流程相连。

\begin{knowledgebox}{应用壁垒三层结构}
\begin{tabularx}{\textwidth}{p{0.20\textwidth}p{0.34\textwidth}X}
\toprule
层级 & 形成方式 & 被复制风险 \\
\midrule
体验壁垒 & 更好的设计师交互和默认流程 & 容易被大模型产品追上。 \\
工作流壁垒 & 项目、团队、版本、资产沉淀 & 迁移成本较高。 \\
生态壁垒 & 工具、素材、品牌和交付系统连接 & 最难复制，但建设最慢。 \\
\bottomrule
\end{tabularx}
\end{knowledgebox}

\begin{warningbox}{垂直应用最怕“薄工作流”}
如果应用只是给模型套一个界面，它会被通用 Agent 很快吃掉。只有当工作流、数据和协作关系真的沉淀下来，垂直应用才有持续位置。
\end{warningbox}

\subsection{本章小结}

Lovart 的长期挑战，是把模型能力窗口转成设计师工作流壁垒。通用 Agent 会收敛入口，但垂直 Agent 仍可能在专业场景里有深水区。

